% ==============================================================================
% CAPITOLO 5: VALORE ATTESO, VARIANZA, COVARIANZA E DISUGUAGLIANZE
% ==============================================================================
\chapter{Valore Atteso, Varianza, Covarianza e Disuguaglianze}
\label{cap:valore_atteso_e_varianza}

\begin{abstract}
Il presente capitolo sviluppa la teoria dei momenti e dei parametri sintetici che governano il comportamento probabilistico delle variabili aleatorie. Viene definito rigorosamente l'operatore lineare di speranza matematica $\mathbb{E}[\cdot]$, fornendo la giustificazione analitica della formula del trasferimento (LOTUS) e dimostrando la proprietà fondamentale dell'additività universale. Si analizzano la varianza teorica (con dimostrazione del Teorema di König-Huygens), la deviazione standard, la covarianza e la correlazione teorica $\rho$ intesa come coseno geometrico nello spazio $L^2$. Il capitolo culmina con le dimostrazioni analitiche complete delle disuguaglianze di concentrazione di Markov e di Chebyshev, pietre miliari per la convergenza stocastica e la Legge dei Grandi Numeri.
\end{abstract}

% ==============================================================================
% SEZIONE 5.1: VALORE ATTESO (SPERANZA MATEMATICA)
% ==============================================================================
\section{Valore Atteso (Speranza Matematica) e Principio del Baricentro}
\label{sec:valore_atteso_teoria}

Se la funzione di probabilità (PMF o PDF) racchiude l'intera descrizione analitica di una variabile aleatoria, per scopi ingegneristici e decisionali è primario sintetizzare tale distribuzione attraverso parametri numerici scalari detti \textbf{momenti}.

Il momento del primo ordine più importante è il **valore atteso** (o speranza matematica / media teorica), che rappresenta la media probabilistica a lungo termine dei valori assunti dalla variabile.

\begin{definizione}{Valore Atteso (Speranza Matematica)}{def_valore_atteso}
Sia $(\Omega, \mathcal{F}, \mathbb{P})$ uno spazio di probabilità e sia $X$ una variabile aleatoria. Il \textbf{valore atteso} di $X$, denotato con $\mathbb{E}[X]$ o $\mu_X$, è definito come:
\begin{itemize}
    \item \textbf{Nel caso Discreto:}
    \begin{equation}
        \mathbb{E}[X] \coloneqq \sum_{x \in \operatorname{Im}(X)} x \cdot p_X(x) = \sum_{x} x \cdot \mathbb{P}(X = x)
    \end{equation}
    a condizione che la serie converga assolutamente: $\sum_{x} |x| p_X(x) < +\infty$.
    \item \textbf{Nel caso Continuo:}
    \begin{equation}
        \mathbb{E}[X] \coloneqq \int_{-\infty}^{+\infty} x \cdot f_X(x) \, dx
    \end{equation}
    a condizione che l'integrale improprio converga assolutamente: $\int_{-\infty}^{+\infty} |x| f_X(x) \, dx < +\infty$.
\end{itemize}
\end{definizione}

\begin{figure}[htbp]
    \centering
    \begin{tikzpicture}[scale=0.85, >=Stealth]
    % Asse x
    \draw[thick, navyblue, ->] (-0.5,0) -- (7.5,0) node[right, font=\footnotesize] {$x$};

    % Blocchi di probabilità (masse)
    \draw[fill=coolblue!30, draw=navyblue, thick] (1,0) rectangle (1.8,1.2);
    \node[font=\tiny\bfseries] at (1.4,0.6) {$p(x_1)$};
    \node[below, font=\footnotesize] at (1.4,0) {$x_1$};

    \draw[fill=coolblue!45, draw=navyblue, thick] (3.2,0) rectangle (4.2,2.0);
    \node[font=\tiny\bfseries] at (3.7,1.0) {$p(x_2)$};
    \node[below, font=\footnotesize] at (3.7,0) {$x_2$};

    \draw[fill=coolblue!25, draw=navyblue, thick] (5.5,0) rectangle (6.3,0.8);
    \node[font=\tiny\bfseries] at (5.9,0.4) {$p(x_3)$};
    \node[below, font=\footnotesize] at (5.9,0) {$x_3$};

    % Fulcro (Baricentro) a x = E[X] = 3.4
    \draw[fill=warmamber, draw=warmamber!90!black, thick] (3.4,0) -- (3.15,-0.6) -- (3.65,-0.6) -- cycle;
    \node[below, font=\bfseries\small, text=warmamber!90!black] at (3.4,-0.65) {$\mathbb{E}[X]$ (Fulcro di Equilibrio)};
    \draw[dashed, warmamber!90!black, thick] (3.4,0) -- (3.4,2.5);

    \node[above, font=\footnotesize\itshape, text=navyblue] at (3.4,2.6) {$\sum (x_i - \mathbb{E}[X]) p(x_i) = 0$};
\end{tikzpicture}

    \caption{Interpretazione fisica e meccanica del valore atteso come baricentro (centro di massa) della distribuzione di probabilità. Il punto $\mu = \mathbb{E}[X]$ rappresenta il fulcro su cui l'asse delle ascisse è in perfetto equilibrio statico rotazionale: $\sum (x_i - \mu)p(x_i) = 0$ oppure $\int (x - \mu)f(x)dx = 0$.}
    \label{fig:baricentro_valore_atteso}
\end{figure}

\subsection{La Formula del Trasferimento (LOTUS)}

Spesso non siamo interessati al valore atteso di $X$, bensì a quello di una sua trasformazione non lineare $Y = g(X)$ (ad esempio, il costo di gestione $C = a X^2 + b$). In linea di principio, occorrerebbe prima determinare la complessa densità $f_Y(y)$ e poi calcolare $\int y f_Y(y) dy$. 

Il celebre Teorema del Trasferimento dimostra che questo passaggio intermedio è superfluo.

\begin{teorema}{Formula del Trasferimento (LOTUS - Law of the Unconscious Statistician)}{lotus}
Sia $X$ una v.a. e sia $g: \mathbb{R} \to \mathbb{R}$ una funzione misurabile. Allora $\mathbb{E}[g(X)]$ può essere calcolato direttamente integrando $g(x)$ rispetto alla misura originaria di $X$:
\begin{align}
    \text{Nel discreto:} \quad & \mathbb{E}[g(X)] = \sum_{x \in \operatorname{Im}(X)} g(x) \cdot p_X(x) \\
    \text{Nel continuo:} \quad & \mathbb{E}[g(X)] = \int_{-\infty}^{+\infty} g(x) \cdot f_X(x) \, dx
\end{align}
Nel caso bivariato per una funzione $h(X, Y)$:
\begin{equation}
    \mathbb{E}[h(X, Y)] = \iint_{\mathbb{R}^2} h(x, y) f_{X,Y}(x, y) \, dx \, dy
\end{equation}
\end{teorema}

\subsection{Linearità Universale del Valore Atteso}

\begin{teorema}{Linearità dell'Operatore Valore Atteso}{linearita_valore_atteso}
Siano $X$ e $Y$ due variabili aleatorie con valore atteso finito e siano $a, b \in \mathbb{R}$. Allora:
\begin{enumerate}
    \item \textbf{Omogeneità e Traslazione:} $\mathbb{E}[aX + b] = a\mathbb{E}[X] + b$.
    \item \textbf{Additività Universale:}
    \begin{equation}
        \mathbb{E}[X + Y] = \mathbb{E}[X] + \mathbb{E}[Y]
    \end{equation}
    \textbf{Attenzione fondamentale:} L'additività del valore atteso è una proprietà algebrica universale dell'integrale e \textbf{vale sempre, sia che $X$ e $Y$ siano indipendenti sia che siano fortemente dipendenti!}
\end{enumerate}
\end{teorema}

\begin{proposizione}{Valore Atteso del Prodotto di Variabili Indipendenti}{prodotto_indipendenti}
Se $X$ e $Y$ sono variabili aleatorie \textbf{stocasticamente indipendenti}, allora:
\begin{equation}
    \mathbb{E}[X \cdot Y] = \mathbb{E}[X] \cdot \mathbb{E}[Y]
\end{equation}
\end{proposizione}

% ==============================================================================
% SEZIONE 5.2: VARIANZA E MOMENTI CENTRATI
% ==============================================================================
\section{Varianza Teorica e Deviazione Standard}
\label{sec:varianza_teorica}

La varianza misura l'entità media della dispersione quadratica della variabile aleatoria attorno al suo valore atteso $\mu$.

\begin{definizione}{Varianza e Deviazione Standard}{def_varianza_teorica}
Sia $X$ una variabile aleatoria con media $\mu = \mathbb{E}[X]$. Si definisce \textbf{varianza} di $X$:
\begin{equation}
    \operatorname{Var}(X) = \sigma^2 \coloneqq \mathbb{E}\left[(X - \mu)^2\right] = \mathbb{E}\left[(X - \mathbb{E}[X])^2\right]
\end{equation}
La \textbf{deviazione standard} (o scarto quadratico medio) è $\operatorname{SD}(X) = \sigma = \sqrt{\operatorname{Var}(X)}$.
\end{definizione}

\begin{teorema}{Formula di Calcolo di König-Huygens per la Varianza}{thm_konig_huygens_teorico}
La varianza ammette la forma computazionale fondamentale:
\begin{equation}
    \operatorname{Var}(X) = \mathbb{E}[X^2] - (\mathbb{E}[X])^2 = \mathbb{E}[X^2] - \mu^2
\end{equation}
\end{teorema}

\begin{dimostrazione}
Sviluppando il quadrato all'interno del valore atteso e applicando la linearità:
\begin{align}
    \operatorname{Var}(X) &= \mathbb{E}\left[X^2 - 2\mu X + \mu^2\right] \\
    &= \mathbb{E}[X^2] - 2\mu \mathbb{E}[X] + \mathbb{E}[\mu^2] \\
    &= \mathbb{E}[X^2] - 2\mu(\mu) + \mu^2 = \mathbb{E}[X^2] - 2\mu^2 + \mu^2 = \mathbb{E}[X^2] - \mu^2 \qquad \blacksquare
\end{align}
\end{dimostrazione}

\begin{proprieta}{Proprietà della Varianza}{prop_varianza}
Per ogni costante $a, b \in \mathbb{R}$:
\begin{enumerate}
    \item $\operatorname{Var}(X) \ge 0$, con $\operatorname{Var}(X) = 0 \iff X = c$ quasi certamente (costante deterministica).
    \item $\operatorname{Var}(aX + b) = a^2 \operatorname{Var}(X)$ (la varianza è insensibile a traslazioni costanti).
    \item $\operatorname{SD}(aX + b) = |a| \operatorname{SD}(X)$.
\end{enumerate}
\end{proprieta}

% ==============================================================================
% SEZIONE 5.3: COVARIANZA E IDENTITÀ DI BIENAYMÉ
% ==============================================================================
\section{Covarianza, Correlazione e Varianza della Somma}
\label{sec:covarianza_teorica}

\begin{definizione}{Covarianza e Correlazione Teorica}{def_covarianza_teorica}
Dati due v.a. $X$ e $Y$ con medie $\mu_X, \mu_Y$:
\begin{enumerate}
    \item La \textbf{covarianza teorica} è definita da:
    \begin{equation}
        \operatorname{Cov}(X, Y) \coloneqq \mathbb{E}[(X - \mu_X)(Y - \mu_Y)] = \mathbb{E}[XY] - \mathbb{E}[X]\mathbb{E}[Y]
    \end{equation}
    \item Il \textbf{coefficiente di correlazione di Pearson} è:
    \begin{equation}
        \rho(X, Y) = \frac{\operatorname{Cov}(X, Y)}{\sqrt{\operatorname{Var}(X)\operatorname{Var}(Y)}} = \frac{\operatorname{Cov}(X, Y)}{\sigma_X \sigma_Y} \in [-1, 1]
    \end{equation}
\end{enumerate}
\end{definizione}

\begin{teorema}{Varianza della Somma e Identità di Bienaymé}{thm_varianza_somma}
Per qualsiasi coppia di variabili aleatorie $X, Y$:
\begin{equation}
    \operatorname{Var}(X + Y) = \operatorname{Var}(X) + \operatorname{Var}(Y) + 2\operatorname{Cov}(X, Y)
\end{equation}
Più in generale, per una combinazione lineare di $n$ variabili:
\begin{equation}
    \operatorname{Var}\left(\sum_{i=1}^n a_i X_i\right) = \sum_{i=1}^n a_i^2 \operatorname{Var}(X_i) + 2 \sum_{1 \le i < j \le n} a_i a_j \operatorname{Cov}(X_i, X_j)
\end{equation}
\textbf{Identità di Bienaymé:} Se le variabili $\{X_i\}_{i=1}^n$ sono \textbf{a due a due incorrelate} (ed in particolare se sono stocasticamente indipendenti, per cui $\operatorname{Cov}(X_i, X_j) = 0$), allora la varianza della somma è esattamente la somma delle varianze:
\begin{equation}
    \operatorname{Var}\left(\sum_{i=1}^n X_i\right) = \sum_{i=1}^n \operatorname{Var}(X_i)
\end{equation}
\end{teorema}

% ==============================================================================
% SEZIONE 5.4: DISUGUAGLIANZE DI CONCENTRAZIONE (MARKOV E CHEBYSHEV)
% ==============================================================================
\section{Disuguaglianze Fondamentali di Concentrazione: Markov e Chebyshev}
\label{sec:disuguaglianze_concentrazione}

In ingegneria è frequente dover stimare una probabilità di guasto o di deviazione estrema conoscendo unicamente la media $\mu$ e la deviazione standard $\sigma$ del sistema, senza conoscerne la densità esatta. Le disuguaglianze di concentrazione forniscono limiti superiori universali rigorosi validi per qualunque distribuzione.

\begin{figure}[htbp]
    \centering
    \begin{tikzpicture}[scale=0.85, >=Stealth]
    \draw[->, thick, navyblue] (-0.5,0) -- (7.5,0) node[right, font=\footnotesize] {$x$};
    \draw[->, thick, navyblue] (0,-0.3) -- (0,3.5) node[above, font=\footnotesize] {$f_X(x)$};

    % PDF
    \draw[domain=0.5:6.5, smooth, variable=\x, thick, navyblue]
        plot ({\x}, {3.0 * exp(-(\x-3.5)*(\x-3.5)/1.8)});

    % Centro mu
    \draw[dashed, navyblue, thick] (3.5,0) -- (3.5,3.0);
    \node[below, font=\small\bfseries] at (3.5,0) {$\mu$};

    % Bande mu - k*sigma e mu + k*sigma
    \draw[thick, crimsonred] (1.6,0) -- (1.6,1.5);
    \draw[thick, crimsonred] (5.4,0) -- (5.4,1.5);
    \node[below, font=\footnotesize, text=crimsonred] at (1.6,0) {$\mu - k\sigma$};
    \node[below, font=\footnotesize, text=crimsonred] at (5.4,0) {$\mu + k\sigma$};

    % Area interna
    \fill[forestgreen!20, domain=1.6:5.4, variable=\x]
        (1.6, 0) --
        plot ({\x}, {3.0 * exp(-(\x-3.5)*(\x-3.5)/1.8)}) --
        (5.4, 0) -- cycle;
    \node[font=\footnotesize\bfseries, text=forestgreen!80!black] at (3.5,1.0) {
        $\mathbb{P}(|X-\mu| < k\sigma) \ge 1 - \frac{1}{k^2}$
    };

    % Code esterne
    \fill[crimsonred!25, domain=0.5:1.6, variable=\x]
        (0.5,0) -- plot ({\x}, {3.0 * exp(-(\x-3.5)*(\x-3.5)/1.8)}) -- (1.6,0) -- cycle;
    \fill[crimsonred!25, domain=5.4:6.5, variable=\x]
        (5.4,0) -- plot ({\x}, {3.0 * exp(-(\x-3.5)*(\x-3.5)/1.8)}) -- (6.5,0) -- cycle;

    \node[above, font=\tiny\bfseries, text=crimsonred] at (1.0, 0.6) {Coda sx};
    \node[above, font=\tiny\bfseries, text=crimsonred] at (6.0, 0.6) {Coda dx};
    \node[above, font=\tiny\bfseries, text=crimsonred] at (3.5, -1.2) {Probabilità totale nelle due code $\le \frac{1}{k^2}$};
\end{tikzpicture}

    \caption{Visualizzazione geometrica della Disuguaglianza di Chebyshev. La probabilità che la variabile $X$ si discosti dalla media $\mu$ per più di $k$ deviazioni standard ($|X - \mu| \ge k\sigma$) è limitata superiormente dall'area delle code esterne pari a $1/k^2$, a prescindere dalla forma specifica della distribuzione.}
    \label{fig:disuguaglianza_chebyshev}
\end{figure}

\begin{teorema}{Disuguaglianza di Markov}{thm_markov}
Sia $Y$ una variabile aleatoria \textbf{non negativa} ($Y \ge 0$ q.c.) con valore atteso finito $\mathbb{E}[Y]$. Allora per ogni costante $a > 0$:
\begin{equation}
    \mathbb{P}(Y \ge a) \le \frac{\mathbb{E}[Y]}{a}
\end{equation}
\end{teorema}

\begin{dimostrazione}
Consideriamo la funzione indicatrice $\IndGraffe{Y \ge a}$ (pari a $1$ se $Y \ge a$, $0$ se $Y < a$). Poiché $Y \ge 0$, vale la disuguaglianza puntuale:
\begin{equation}
    a \cdot \IndGraffe{Y \ge a} \le Y \qquad (\forall \omega \in \Omega)
\end{equation}
Infatti:
\begin{itemize}
    \item Se $Y(\omega) < a$, il membro sinistro vale $a \cdot 0 = 0 \le Y(\omega)$;
    \item Se $Y(\omega) \ge a$, il membro sinistro vale $a \cdot 1 = a \le Y(\omega)$.
\end{itemize}
Applicando la monotonia del valore atteso ad ambo i membri:
\begin{equation}
    \mathbb{E}\left[a \cdot \IndGraffe{Y \ge a}\right] \le \mathbb{E}[Y] \iff a \cdot \mathbb{E}\left[\IndGraffe{Y \ge a}\right] \le \mathbb{E}[Y]
\end{equation}
Poiché per definizione $\mathbb{E}\left[\IndGraffe{Y \ge a}\right] = 1 \cdot \mathbb{P}(Y \ge a) + 0 \cdot \mathbb{P}(Y < a) = \mathbb{P}(Y \ge a)$, dividendo per $a > 0$ otteniamo:
\begin{equation}
    \mathbb{P}(Y \ge a) \le \frac{\mathbb{E}[Y]}{a} \qquad \blacksquare
\end{equation}
\end{dimostrazione}

\begin{teorema}{Disuguaglianza di Chebyshev (Bienaymé-Chebyshev)}{thm_chebyshev}
Sia $X$ una variabile aleatoria con media $\mu = \mathbb{E}[X]$ e varianza $\sigma^2 = \operatorname{Var}(X) < +\infty$. Allora per ogni soglia $\varepsilon > 0$:
\begin{equation}
    \mathbb{P}(|X - \mu| \ge \varepsilon) \le \frac{\operatorname{Var}(X)}{\varepsilon^2} = \frac{\sigma^2}{\varepsilon^2}
\end{equation}
Esprimendo la soglia come multiplo $k > 0$ della deviazione standard ($\varepsilon = k\sigma$):
\begin{equation}
    \mathbb{P}(|X - \mu| \ge k\sigma) \le \frac{1}{k^2}, \qquad \text{o equivalentemente} \qquad \mathbb{P}(|X - \mu| < k\sigma) \ge 1 - \frac{1}{k^2}
\end{equation}
\end{teorema}

\begin{dimostrazione}
Definiamo la variabile non negativa $Y = (X - \mu)^2 \ge 0$. Notiamo che la condizione $|X - \mu| \ge \varepsilon$ equivale a $(X - \mu)^2 \ge \varepsilon^2$.
Applicando la Disuguaglianza di Markov a $Y$ con soglia $a = \varepsilon^2$:
\begin{equation}
    \mathbb{P}(|X - \mu| \ge \varepsilon) = \mathbb{P}\left((X - \mu)^2 \ge \varepsilon^2\right) \le \frac{\mathbb{E}\left[(X - \mu)^2\right]}{\varepsilon^2} = \frac{\operatorname{Var}(X)}{\varepsilon^2} = \frac{\sigma^2}{\varepsilon^2} \qquad \blacksquare
\end{equation}
\end{dimostrazione}

\paragraph{Esempio di Utilizzo Ingegneristico:}
Senza conoscere nulla sulla legge di probabilità di un macchinario:
\begin{itemize}
    \item Per $k = 2$: la probabilità di trovarsi fuori dall'intervallo $[\mu - 2\sigma, \mu + 2\sigma]$ non può superare $\frac{1}{2^2} = 25\%$;
    \item Per $k = 3$: la probabilità di deviazione oltre $3\sigma$ non può superare $\frac{1}{3^2} = 11.11\%$.
\end{itemize}
